\section{Placement: stations, seats, and the resolver}\label{sec:placement}

\Cref{sec:parts} described the rocket as an ownership tree of parts and the machinery that
folds that tree into mass, center of mass (CM), and inertia. This section describes the
other half of the model layer's geometry story: how the tree knows \emph{where} each part
sits. The answer is deliberately indirect. QtRocket never stores a coordinate for any part.
What it stores, on every parent$\to$child edge, is a small value called a
\code{StationLink} --- a statement of \emph{physical intent} (``this rim seats against that
rim, inserted this deep'') --- and every absolute position in the rocket is recomputed from
those statements by exactly one resolver function whenever the structure changes
\src{model/parts/Placement.h}{147}. The section walks the design from philosophy to
enforcement: the rules, the data model (stations, seats, links), the resolver arithmetic
term by term, the diagnostics that refuse impossible geometry, the narrow authoring
surface that exists today, and the invariance suite that pins the system's observable
behavior.

\subsection{Philosophy: placement is intent, never coordinates}\label{sec:placement-philosophy}

The placement system was designed in its own whitepaper
(\srcf{docs/PartPlacementDesignLatex/}) and implemented essentially as designed; where the
document and the shipped code disagree, this section reports the code and says so. The
design opens with three rules stated before any type is introduced, because they constrain
every later decision.

\begin{designrule}[Three non-negotiable rules]
(1) \emph{Placement is derived}: absolute position is recomputed by exactly one resolver
--- neither an absolute coordinate nor a resolved pose is ever stored as truth.
(2) \emph{The stored link is physical intent}: a feature-to-feature relationship between
two parts, never an opaque offset; the center of mass falls out of geometry as a
consequence. (3) \emph{The common case costs zero authored numbers}: a child attached with
no link simply abuts its parent
\src{docs/PartPlacementDesignLatex/sections/02-philosophy.tex}{21}.
\end{designrule}

The rules are a reaction to an observed failure. The representation they replaced stored
one \code{Vector3} per child, defined as the child's CM relative to the parent's CM. A CM
is a derived quantity --- a solid nose cone's sits a quarter length from the base, not at
mid-length --- so every offset was a hand calculation, and the value's two consumers
quietly disagreed about what it meant: the simulator applied it CM-to-CM as documented,
the visualizer as a geometric-center-to-center translation. The two reconstructions of
the same file diverged by exactly the cone's CM offset, and a committed design shipped
with a coupler silently interpenetrating the nose cone --- no check fired, and a
self-intersecting solid was fed to the integrator
\src{docs/PartPlacementDesignLatex/sections/01-introduction.tex}{65}. The link recorded a
derived quantity where the author actually knew a geometric fact --- two rims touch. The
replacement stores the geometric fact.

Storing intent also leaves nothing to go stale. Persisting the resolved transform
alongside the intent was rejected in the design document as the same staleness bug
reflected: edit a length and the stored transform is silently wrong
\src{docs/PartPlacementDesignLatex/sections/02-philosophy.tex}{43}. Resolved poses exist
in the shipped code only as a cache gated on a structural dirty flag
(\cref{sec:placement-resolver}) --- an optimization that can never be wrong, not a second
copy of the truth that can.

\begin{designrule}[One resolver, two consumers, identical numbers]
``The single authority for absolute placement: intent (StationLink) in, geometry (Pose)
out. Both the simulator and the visualizer consume its output, so the two cannot
disagree'' \src{model/parts/Placement.h}{148}. The mass-properties walk reads the
resolver's cached output when it folds the tree \src{model/parts/Part.cpp}{199}, and the
visualizer's mesh walk reads the same output when it draws it
\src{visualizer/RocketMesh.cpp}{434} --- so the number the integrator uses to place a
part's mass is, by construction, the number the renderer uses to draw it. The legacy
divergence is structurally impossible.
\end{designrule}

One paragraph of contrast situates the choice. OpenRocket --- the comparison baseline of
\cref{app:or-architecture} --- positions each component with a signed axial offset plus an
\code{AxialMethod} enumerator (top, middle, bottom, absolute) selecting which datum of the
parent the offset is measured from. That scheme also re-derives absolute coordinates from
live parent geometry, so it is not a staleness hazard; the difference is semantic. The
stored number is a coordinate displacement whose meaning is switched by a flag, and the
physical relationship it encodes --- \emph{is} this coupler inserted into that tube, or
merely near it? --- is not represented and cannot be checked. QtRocket's
\code{StationLink} stores the mating itself: which stations of parent and child sit
together, and in what radial relationship. That extra semantic content is what gives the
diagnostics layer of \cref{sec:placement-diagnostics} something to enforce, and it is why
QtRocket can afford to fail closed where OpenRocket's tradition is to warn and continue.

\subsection{The data model: stations, links, and seats}\label{sec:placement-datamodel}

\begin{physbox}[Stations, fore and aft]
Aerospace drafting locates features on a vehicle by \emph{station}: a cross-section plane
at a stated distance from a reference datum --- ``fuselage station 172'' names a plane,
not a part. \emph{Fore} means toward the nose, \emph{aft} toward the tail. A rocket
airframe is, to good approximation, a stack of axisymmetric bodies, so a station plus the
radius profile there --- outer skin and, for hollow parts, inner bore --- captures
everything about the cross-section that matters for fit. QtRocket's placement model is
this drafting idea made executable: parts are mated station-to-station, and the radius
profiles at the mated stations decide whether the mating is physically possible.
\end{physbox}

Every part lives in one shared local frame, established in \cref{sec:parts}: \zfwd, the
part's own origin at its \emph{fore} plane, the body occupying $z \in [-L, 0]$
\src{model/parts/Part.h}{130}. Stations are expressed as fractions: \code{station01}
$\in [0,1]$ with $0$ the aft plane and $1$ the fore plane, mapped to a local coordinate by

\begin{equation}\label{eq:placement-station}
z(\mathtt{station01}) \;=\; (\mathtt{station01} - 1)\cdot L ,
\end{equation}

where $L$ is the part's \emph{live} length, re-read from \code{getLength()} on every
resolve \src{model/parts/Part.cpp}{305}. \code{Part::stationAt} clamps the fraction to
$[0,1]$, evaluates \cref{eq:placement-station}, and returns a \code{Station} landmark ---
$\{z, r_{\mathrm{outer}}, r_{\mathrm{inner}}\}$, the coordinate plus the radius profile
read at it --- which is always derived, never stored \src{model/parts/Placement.h}{62}.
Fractions of live length are the anti-staleness property in miniature: lengthen a tube and
every seam expressed against it moves to follow, because no stored number baked the old
length in (pinned by test: \src{model/tests/ResolverSweepTests.cpp}{168}). One honest
caveat: at the pinned commit \code{254cd41}, concrete parts are geometrically immutable
after construction --- no length or radius setters exist \src{model/parts/BodyTube.h}{69}
--- so the only structural edits are attach and detach, and the live-length machinery is
future-proofing for an editing path not yet built (\cref{sec:incomplete}).

The stored value on each edge is the \code{StationLink}
(\cref{lst:placement-stationlink}): a fractional station on each of the two mated parts, a
gap in meters, a seat kind, and a per-seam quaternion reserved for six-degree-of-freedom
(6-DOF) work --- identity today, so a stored link is bit-identical to a pure translation
\src{model/parts/Placement.h}{44}.

\begin{listing}[htbp]
\begin{minted}{cpp}
struct StationLink
{
   double     parentStation01{0.0};             ///< 0 = aft plane, 1 = fore plane, on the PARENT
   double     childStation01{1.0};              ///< 0 = aft plane, 1 = fore plane, on the CHILD
   double     gap{0.0};                         ///< meters; meaning per SeatKind
   SeatKind   seat{SeatKind::Abut};             ///< selects the radial check and the gap sign
   Quaternion childRot{Quaternion::Identity()}; ///< (x,y,z,w) seam orientation; identity in 3-DOF
};
\end{minted}
\caption{The only stored spatial relationship in the model \srccap{model/parts/Placement.h}{49}.
The defaults are load-bearing: a default-constructed link means ``child fore plane abuts
parent aft plane,'' the natural nose$\to$body$\to$fins build order, so the common case
costs zero authored numbers.}\label{lst:placement-stationlink}
\end{listing}

The \code{SeatKind} enumeration is closed at three values by design: a rigid axisymmetric
stack admits exactly three radial relationships \src{model/parts/Placement.h}{32}, and
leaving the enum open ``would only invite half-defined cases''
\src{docs/PartPlacementDesignLatex/sections/03-datamodel.tex}{126}. The seat does double duty:
it names the radial compatibility semantics of the seam, and it fixes the \emph{sign} of
the gap --- the direction in which the gap displaces the child
(\cref{tab:placement-seats}). Encoding direction in the seat, rather than asking the author
to sign the number, is what lets an insertion depth of $0.04\m$ and a forward standoff of
$0.04\m$ both be written as a plain positive number. The non-negativity of a
\code{NestInBore} depth is a documented contract of the authoring helpers, not a runtime
check \src{model/parts/Placement.h}{133}; an out-of-contract value surfaces only if the
resulting geometry trips the envelope sweep.

\begin{table}[htbp]
\centering\small
\begin{tabularx}{\linewidth}{@{}l>{\raggedright\arraybackslash}X>{\raggedright\arraybackslash}X@{}}
\toprule
Seat kind & Radial compatibility semantics & Gap meaning and sign \\
\midrule
\code{Abut} & Rim to rim: the two outer radii at the mated stations must match,
  $|p.r_{\mathrm{outer}} - c.r_{\mathrm{outer}}| \le \mathrm{tol}$
  \src{model/parts/Placement.cpp}{66} & Forward standoff along $+z$; flush at zero. The
  default seat. \\
\addlinespace
\code{NestInBore} & Child outer diameter (OD) fits inside the parent's inner diameter (ID):
  $c.r_{\mathrm{outer}} \le p.r_{\mathrm{inner}} + \mathrm{tol}$
  \src{model/parts/Placement.cpp}{72} & Insertion \emph{depth}, driving the child aft
  ($-z$); the one sign-flipping seat. \\
\addlinespace
\code{OnSurface} & Child seats radially on the parent's outer wall: child radius equals
  the parent's outer radius at the seat station \src{model/parts/Placement.cpp}{69} &
  Axial standoff along the wall. \\
\bottomrule
\end{tabularx}
\caption{The three seat kinds \srccap{model/parts/Placement.h}{32}. Each seat selects both
a radial check between the two mated \code{Station} landmarks and the direction the gap
displaces the child. The radial semantics are implemented in \code{radialSeamCheck}; the
enforcement status of that check is examined in \cref{sec:placement-diagnostics}.}
\label{tab:placement-seats}
\end{table}

Three snap-operator verbs provide sugar over the brace-initializer form:
\cppinline|abut(gap)|, \cppinline|nestInBore(depth)|, and
\cppinline|seatOnWall(parentStation01)| each produce a \code{StationLink} by value
\src{model/parts/Placement.h}{127}. Their defining property is what they do \emph{not} do:
no verb reads a length, computes a station, stores a pose, or mutates a part
\src{model/parts/Placement.h}{122}. Verbs name intent; only the resolver resolves it ---
so a verb called before its operands are sized, or a link reused across designs, cannot
bake in a stale coordinate.

\tikzfig{placement-stationlink}{Anatomy of a \code{StationLink}, shown on the
\code{xl75\_multi} stack that exercises all three seat kinds. Each parent$\to$child edge
stores fractional stations on both parts, a seat kind, and a seat-signed gap; the resolver
turns each link into a fore-plane pose in the nose-tip frame. The coupler's link
(\code{NestInBore}, depth $0.04\m$ into the body's fore bore) resolves to origin
$z=-0.26\m$; the fins' link (\code{OnSurface} at parent station $0.06$) seats them on the
tube wall near its aft end. No stored number is an absolute coordinate.
\srccap{model/parts/Placement.h}{49}}{fig:placement-stationlink}

\subsection{The resolver, term by term}\label{sec:placement-resolver}

The resolver's output frame is fixed by one convention: \code{resolvePlacements(root,
Pose\{\})} plants the root part's fore plane --- for a conventional rocket, the nose tip
--- at the origin, and the stack grows aft into $-z$
\src{model/parts/Placement.cpp}{38}. Every derived quantity downstream inherits this
\emph{nose-tip datum}: resolved part origins, the composite CM, and the composite center
of pressure are all coordinates on the same tip-anchored axis, which is what makes the
static margin of \cref{sec:aero} a literal subtraction. The datum is a fixed, recognizable
geometric feature rather than anything mass-dependent --- a deliberate contrast with the
retired CM-relative convention.

Placing one child is a pure function of the parent's resolved pose, the two parts'
geometry, and the link (\cref{lst:placement-resolver}).

\begin{listing}[htbp]
\begin{minted}{cpp}
   const Station p = parent.stationAt(link.parentStation01);  // landmark on the parent
   const Station c = child.stationAt(link.childStation01);    // landmark on the child

   // signedGap encodes seat direction in +z = forward: NestInBore inserts aft (-z, gap is depth),
   // Abut/OnSurface stand off forward (+z). The author always writes a non-negative number.
   const double signedGap = (link.seat == SeatKind::NestInBore) ? -link.gap : link.gap;

   // Line the child's station up with the parent's, then displace by the gap. The pose locates the
   // child's fore plane, which sits c.z from its chosen station -- hence the `- c.z`.
   const double childOriginZ = p.z + signedGap - c.z;

   const Pose childInParent{Vector3(0.0, 0.0, childOriginZ), link.childRot};  // coaxial in 3-DOF
   return parentPose.compose(childInParent);
\end{minted}
\caption{The entire resolver core \srccap{model/parts/Placement.cpp}{23}: two station
lookups, one seat-signed gap, one subtraction, one rigid-transform composition. Everything
else in the placement system exists to feed this arithmetic or to check its output.}
\label{lst:placement-resolver}
\end{listing}

\begin{equation}\label{eq:placement-resolver}
\mathtt{childOriginZ} \;=\; p.z \;+\; \mathtt{signedGap} \;-\; c.z,
\qquad
\mathtt{signedGap} \;=\;
\begin{cases}
-\,\mathtt{gap} & \text{\texttt{NestInBore},}\\
+\,\mathtt{gap} & \text{\texttt{Abut}, \texttt{OnSurface}.}
\end{cases}
\end{equation}

Each term earns its place:

\begin{itemize}
\item $p.z$ is the parent-local coordinate of the chosen parent station,
  $(\mathtt{parentStation01}-1)\cdot L_{\mathrm{parent}}$ by
  \cref{eq:placement-station} --- \emph{where on the parent the seam sits}.
\item $\mathtt{signedGap}$ is the authored gap with the seat's direction applied. A
  \code{NestInBore} gap is an insertion depth, so it moves the child aft ($-z$); an
  \code{Abut} or \code{OnSurface} gap is a standoff, so it moves the child forward
  ($+z$) \src{model/parts/Placement.cpp}{28}. The author never signs the number; the
  seat does.
\item $-\,c.z$ converts a seam position into a fore-plane position. The \code{Pose}
  locates the child's \emph{fore plane} --- its local origin --- but the seam was
  expressed against the child's chosen station, which sits $c.z$ from that origin.
  Subtracting $c.z$ re-expresses the one as the other
  \src{model/parts/Placement.cpp}{30}.
\end{itemize}

The resulting child-in-parent pose is $(0, 0, \mathtt{childOriginZ})$ with orientation
\code{link.childRot}: $x = y = 0$ because placement is coaxial in three degrees of freedom
(3-DOF). Radial position is deliberately a \emph{check}, not a coordinate --- the radii
recovered from the two stations feed the diagnostics, but no pose is ever displaced off
axis \src{model/parts/Placement.cpp}{34}. The design document rejects off-axis seating in
3-DOF explicitly: without orientation and a body-frame CM offset it would store a
half-correct placement, the exact defect class the design exists to eliminate, relocated
from the axial axis to the radial one
\src{docs/PartPlacementDesignLatex/sections/04-resolver.tex}{132}. The pose is composed
into the root frame by \code{Pose::compose} --- rotate-then-translate, which under
identity quaternions degenerates bit-identically to vector addition
\src{model/parts/Placement.h}{81}; carrying the quaternion anyway is the zero-schema path
to 6-DOF (\cref{sec:incomplete}).

\code{resolvePlacements} is then a depth-first search (DFS) from the root pose, calling
\code{placeChild} for each \code{(child, link)} pair in attachment order and emitting one
\code{Placed} --- a non-owning part pointer plus its resolved pose --- per part, in
deterministic pre-order \src{model/parts/Placement.cpp}{38}. The determinism is pinned by
test \src{model/tests/ResolverSweepTests.cpp}{151} because the overlap sweep's
reproducibility depends on it. The resolver reads ``no CM, time, or mass''
\src{model/parts/Placement.h}{158}: it is pure geometry, which is what allows it to run on
a completely different cadence from the mass machinery.

Worked numbers make the arithmetic concrete, on the \code{xl75\_multi} stack of
\cref{fig:placement-stationlink}; every value below is pinned in
\src{model/tests/ResolverSweepTests.cpp}{95} and reproduced end-to-end through the design
file reader in \src{tests/PokeThroughRegressionTests.cpp}{70}. The nose (root, length
$0.30\m$) plants at the identity pose: tip at $z=0$, base rim at $-0.30\m$. The body tube
(length $0.90\m$) attaches with the default link --- zero authored numbers --- so
$p.z=-0.30$, $c.z=0$, gap $0$: its fore plane lands at $-0.30\m$ and it spans
$[-1.20,-0.30]$. The fins seat on the wall at parent station $0.06$: $p.z =
(0.06-1)\cdot0.90 = -0.846$ (world $-1.146$), $c.z = -0.10$ (root chord), so the fin
origin resolves to world $-1.046\m$. The interesting case is the coupler (length
$0.08\m$), nested \cppinline|nestInBore(0.04)| in the body's fore bore:
\[
\mathtt{childOriginZ} = \underbrace{0}_{p.z\ \text{(bore mouth)}}
  \underbrace{-\,0.04}_{\mathtt{signedGap}}
  \underbrace{+\,0.08}_{-\,c.z\ \text{(aft plane)}} = +0.04
\]
in the body frame --- world $-0.26\m$, spanning $[-0.34, -0.26]$. Its aft $0.04\m$ is
correctly buried in the bore; its fore $0.04\m$ projects forward past the body rim, a
fact the resolver happily records and \cref{sec:placement-diagnostics} refuses.

Where do the poses land? Each \code{Part} memoizes its own sub-tree's resolution in a
mutable \code{resolvedCache}, rebuilt by \code{ensurePlacementCache} only when the
\emph{structural} dirty flag \code{placementDirty} is set --- by attach and detach, never
by a mass or inertia edit \src{model/parts/Part.h}{264}. The rebuild does two things in
one breath: it re-resolves the poses (this part planted at the local origin) and re-runs
the envelope sweep of \cref{sec:placement-diagnostics}, caching the verdict alongside the
poses \src{model/parts/Part.cpp}{160}. Two in-tree consumers read the cache: the two-pass
composite mass/CM/inertia walk of \cref{sec:parts} \src{model/parts/Part.cpp}{199} and
the composite aerodynamics fold of \cref{sec:aero} \src{model/parts/Part.cpp}{231}. The
visualizer's mesh walk invokes the same resolver on the same tree --- ``consume the same
absolute placement the simulator does,'' its own comment says --- and reads the cached
diagnostics verdict alongside it \src{visualizer/RocketMesh.cpp}{434}.

\begin{keyinsight}[Two gates, two cadences]
The mass-delta gate of \cref{sec:parts} fires every step while a motor burns, because
composite mass moves every step; the placement gate fires only when the rocket is
structurally rebuilt, because geometry is invariant under a burn
\src{model/parts/Part.cpp}{162}. A white-box test proves the split: during a simulated
burn the composite CM shifts every step while a counting part's radius-profile call
count stays frozen --- the resolver and the sweep never re-fire inside the
ordinary-differential-equation (ODE) loop \src{model/tests/DiagnosticsGateTests.cpp}{97}.
\end{keyinsight}

\subsection{Diagnostics: the envelope sweep and the fail-closed gate}\label{sec:placement-diagnostics}

Deriving placement from intent removes the consumers-disagree error class, but not the
other one: a set of links can be individually sensible and still describe a physically
impossible rocket --- the worked coupler above pokes $0.04\m$ out of its bore into the
solid nose. The resolver places it without complaint; placement is pure geometry. The
diagnostics layer is the component that refuses such designs, and refuses them with a
\emph{located} reason rather than a silent miscomputation.

The enforced check is the \emph{envelope sweep}, \code{sweepOverlaps}
\src{model/parts/Placement.cpp}{107}: a global one-dimensional spatial query over the
resolved tree, run from \code{ensurePlacementCache} once per structural change
\src{model/parts/Part.cpp}{160}. Its steps:

\begin{enumerate}
\item \emph{Intervals.} Each resolved part contributes a world axial interval
  $[z_{\mathrm{aft}}, z_{\mathrm{fore}}] = [\mathtt{originZ} - \mathtt{axialLength},\,
  \mathtt{originZ}]$ --- valid because orientation is identity in 3-DOF
  \src{model/parts/Placement.cpp}{95}.
\item \emph{Sort.} Intervals sort by aft station, making host lookup a sweep --- the
  $O(N \log N)$ in the design's name \src{model/parts/Placement.cpp}{151}.
\item \emph{Feature breakpoints.} Each offender is sampled at its own two endpoints plus
  every other interval endpoint within its span \src{model/parts/Placement.cpp}{160} ---
  exact rather than sampled, \emph{given} the premise that radius profiles are
  closed-form and monotone between breakpoints (examined below).
\item \emph{Host selection.} At each breakpoint the host is the smallest-id covering
  interval, excluding the offender itself, its descendants (legitimate nesting), and its
  direct seat parent \src{model/parts/Placement.cpp}{142}. Smallest-id is a
  deterministic, reload-stable tie-break, pinned by test
  \src{model/tests/ResolverSweepTests.cpp}{325}; stated plainly, the check is
  single-host per breakpoint --- a violation against a larger-id covering host goes
  unevaluated at that station.
\item \emph{Capacity test.} The offender's outer radius is tested against the host's
  capacity under the solid-host rule, with a guard exempting an offender on or outside a
  hollow host's own skin --- a modeling artifact (a fin set's body disc over a co-radial
  aft coupler), not a collision \src{model/parts/Placement.cpp}{192}.
\item \emph{Verdict.} The deepest violation per (offender, host) pair is kept;
  \code{SolveResult.ok} holds exactly when the diagnostics list is empty
  \src{model/parts/Placement.cpp}{230}.
\end{enumerate}

\begin{designrule}[The solid-host rule]
What radius does a host \emph{offer} an intruder? For a bored part, its bore; for a solid
part, its outer skin. \cppinline|return isSolid() ? radiusOuterAt(zLocal) : radiusInnerAt(zLocal);|
\src{model/parts/Part.cpp}{314} --- and an offender of radius $r$ fits iff
$r \le \mathtt{innerCapacityAt}(z) + \mathrm{tol}$ \src{model/parts/Part.h}{158}. Without
the rule, a naive OD-versus-bore predicate reads a solid cone's bore as zero everywhere
and flags every neighbor; with it, the same single predicate is simultaneously the
nesting check and the poke-through check, and call sites never branch on solidity.
\end{designrule}

On the worked stack the sweep catches exactly the violation the resolver recorded. Over
$[-0.34, -0.30]$ the only interval covering the coupler is the body's --- its direct seat
parent, excluded from host selection --- so no capacity test fires there. Over
$[-0.30, -0.26]$ the covering envelope that survives exclusion is the solid nose --- a \emph{non-tree
neighbor}: the coupler is linked to the body, not the nose, and the sweep sees the
collision precisely because it does not restrict itself to the ownership tree. At the
breakpoint $z=-0.26$ the nose skin has tapered to $0.0395 \cdot (0.26/0.30) =
0.0342\m$, so the coupler's OD exceeds capacity by $0.0034\m$: the diagnostic names
offender and host by stable id, the world station, and the penetration depth
\src{model/tests/ResolverSweepTests.cpp}{249}. Ids, not names or pointers, are the only
safe handle: part names need not be unique, and a raw pointer would dangle if the tree
were edited between resolve and report \src{model/parts/Placement.h}{99}.

The verdict is consumed fail-closed. \code{computeCompositeAt} --- the single walk behind
the composite mass-property reads --- refuses a failed solve before doing any arithmetic
(\cref{lst:placement-gate}), so both \code{getCompositeCm(t)} and \code{getCompositeI(t)}
throw on a self-intersecting design \src{model/parts/Part.cpp}{148}; the behavior is
pinned directly \src{model/tests/DiagnosticsGateTests.cpp}{66} and end-to-end through the
design-file reader on a frozen poke-through fixture
\src{tests/PokeThroughRegressionTests.cpp}{46}. Exactly one composite reader is
deliberately exempt: \code{getCompositeMass(t)} stays an ungated live sum, because it is
the hot ODE divisor and does no placement work
\src{model/tests/DiagnosticsGateTests.cpp}{81}. The visualizer reads the same cached
verdict through \code{placementDiagnostics()} \src{model/parts/Part.h}{125} and flags the
offending part \src{visualizer/RocketMesh.cpp}{440} --- one verdict, two consumers,
computed where the geometry changes rather than where it is consumed, so guarding the
inner loop costs a single boolean read.

\begin{listing}[htbp]
\begin{minted}{cpp}
   // A self-intersecting design is a hard, located failure, never a silently-wrong mass/inertia. The
   // verdict is cached (computed once per structural resolve), so this is a flag read, not a re-sweep.
   if(!resolvedDiagnostics.ok)
   {
      const std::string detail = resolvedDiagnostics.diagnostics.empty()
                                    ? std::string("self-intersecting geometry")
                                    : resolvedDiagnostics.diagnostics.front().message;
      throw std::runtime_error("Part::computeCompositeAt: placement solve failed -- " + detail);
   }
\end{minted}
\caption{The fail-closed gate at the top of the composite walk
\srccap{model/parts/Part.cpp}{182}: the thrown error carries the first located diagnostic,
so the caller learns who intruded into what, where, and by how much.}
\label{lst:placement-gate}
\end{listing}

\begin{hardfail}[The monotonicity premise fails for the sphere]
The breakpoint scheme's exactness rests on the claim that every radius profile is
``closed-form and monotone between breakpoints'' \src{model/parts/Placement.cpp}{161}.
True for tubes (constant), cones (linear taper), and fin sets (body disc) --- false for
\code{HollowSphere}, whose silhouette $\sqrt{r_o^2 - (z + r_o)^2}$ is zero at both poles
(its own interval endpoints) and maximal at the mid-span equator
\src{model/parts/HollowSphere.cpp}{53}. A sphere nested in a too-narrow tube evaluates
$r \approx 0$ at every breakpoint and passes undetected; no sweep test exercises a sphere
\srcf{model/tests/ResolverSweepTests.cpp}. The placement design document itself listed
``a sphere bulges'' among the monotone profiles --- one of the places the companion
documents disagree with the shipped code, and the Architecture Review's code-verified
finding wins. \fail
\end{hardfail}

Two further layers of the designed diagnostics stack deserve an honest status report. The
\emph{radial seam check} --- Layer 1 in the design document, the per-seam radial
compatibility test of \cref{tab:placement-seats} --- is implemented in full as
\code{radialSeamCheck} \src{model/parts/Placement.cpp}{56} and unit-tested across all
three seats \src{model/tests/ResolverSweepTests.cpp}{188}, but it is \emph{dormant}: the
cached verdict contains Layer-2 sweep output only \src{model/parts/Part.cpp}{170}, and no
production code calls the function. The design document places it ``on every resolve''
\src{docs/PartPlacementDesignLatex/sections/06-diagnostics.tex}{66}; the as-built
resolver does not run it, and its \code{zWorld} field is left parent-local behind the
comment ``the root-frame z is added once a pose is known''
\src{model/parts/Placement.cpp}{87} --- a promise no caller exists to fulfill. The
practical consequence: a pure rim-radius mismatch producing no axial-interval overlap is
not in the cached verdict today. Layer 3 --- degree-of-freedom accounting, a detector
warning when a child is under-constrained under free orientation --- exists only in the
design document \src{docs/PartPlacementDesignLatex/sections/06-diagnostics.tex}{179}; no
corresponding code exists anywhere in \srcf{model/parts/}. Both gaps are inventoried in
\cref{sec:incomplete}.

\subsection{Authoring today: one verb reaches the front ends}\label{sec:placement-authoring}

The placement model's vocabulary is far richer than any interface that speaks it. At the
pinned commit, exactly one placement operation is reachable from a front end: the
command-line interface's (CLI's) \code{addpart} command, which attaches every new part
with \cppinline|model::part::abut(offset.z())| --- the child's fore plane against the
parent's aft plane, with an optional forward standoff taken from a parsed \code{z=} token
\src{cli/Repl.cpp}{872}. The grammar accepts \code{x=} and \code{y=} too, but they are
parsed and discarded: placement is coaxial in 3-DOF, and only the axial component reaches
the link \src{cli/Repl.cpp}{189}. The graphical interface authors no structure at all at
this commit (\cref{sec:interfaces}), so the CLI's single verb \emph{is} the interactive
authoring surface.

Everything else arrives by other roads. The full seat vocabulary ---
\code{NestInBore}, \code{OnSurface}, non-default stations, arbitrary gaps --- is reachable
from a design file, where the serializer reads a \code{<link>} element's seat, stations,
and gap \src{model/DesignSerializer.cpp}{186} and rejects an unknown seat string
fail-closed \src{model/DesignSerializer.cpp}{184} (\cref{sec:persistence}), and from C++,
where the one production author of a link with non-default stations is the motor attachment:
\code{RocketModel::setMotorModel} joins the airframe root's CM station to the motor's CM
station with an \code{Abut} link whose gap is a motor offset that is zero today
\src{model/RocketModel.cpp}{124}. Nothing in production sets the reserved \code{childRot}
to a non-identity value. A richer authoring surface --- interactive seat selection,
station editing, the geometry mutation that would finally exercise the live-length
machinery --- is future work the data model was expressly shaped for;
\cref{sec:incomplete} carries the inventory.

\subsection{The invariance gate: the system's contract}\label{sec:placement-invariance}

The placement system replaced a load-bearing representation underneath a working
simulator, and its correctness argument is embodied in a test suite that predates the
implementation: \srcf{tests/PlacementInvarianceTests.cpp}, written as Phase~0 of the
migration before any production type changed \src{tests/PlacementInvarianceTests.cpp}{1}.
The gate loads the 24 committed design fixtures \src{tests/PlacementInvarianceTests.cpp}{66}
and compares, per design, the quantities the integrator consumes against an immutable
committed baseline generated by the \emph{legacy} CM-to-CM reader before the refactor:

\begin{itemize}
\item composite mass and composite inertia about the CM --- datum-independent quantities
  --- match the frozen baseline to a relative $10^{-9}$, a tolerance chosen to admit
  unit-in-the-last-place drift and nothing else
  \src{tests/PlacementInvarianceTests.cpp}{147};
\item the composite CM matches after one known, deliberate correction: the datum change
  from the root's own CM to the nose tip, re-expressed by the single fixed offset of the
  root's local CM station \src{tests/PlacementInvarianceTests.cpp}{286};
\item every part's resolved axial station matches under the same single shift
  \src{tests/PlacementInvarianceTests.cpp}{299};
\item the static margin is deliberately \emph{not} compared: the legacy center of
  pressure was contaminated by part CMs, and correcting it was part of the point
  \src{tests/PlacementInvarianceTests.cpp}{318};
\item a save--reload round trip through the current design writer reproduces the same
  snapshot, pinning idempotence of the intent serialization
  \src{tests/PlacementInvarianceTests.cpp}{358}.
\end{itemize}

Two hand-computed pins guard the specific defects the migration corrected: the solid nose
cone's local CM at $-0.225\m$ for $L = 0.30\m$ --- three quarters of the length aft of
the tip, computed by hand independently of both code paths, isolating the cone's
historical reversed-frame sign \src{tests/PlacementInvarianceTests.cpp}{445} --- and the
fin set's axial CM reference, worth half a root chord, guarding the end-versus-mid-chord
confusion \src{tests/PlacementInvarianceTests.cpp}{458}.

\begin{keyinsight}[Invariance is defined by the incumbent, not the spec]
The baseline is the legacy reader's literal output --- including its mis-placements ---
not a re-derivation from each design's intended geometry, because a re-derivation could
silently share whatever bug the legacy path contained. The file's own header declares it
immutable ground truth that must never be regenerated to make a failing migration pass
\src{tests/PlacementInvarianceTests.cpp}{30}. Correcting a design's actual geometry (the
interpenetrating coupler) was done separately and visibly: the fixed corpus resolves
clean, and the original defect survives as a frozen regression fixture that must keep
failing \src{tests/PokeThroughRegressionTests.cpp}{46}.
\end{keyinsight}

The gate is the placement system's contract in the strongest sense available to a
refactor: every observable the flight path consumes is pinned across the representation
change, every deliberate deviation (the tip datum, the corrected CP) is named and
accounted for, and the one behavior that was supposed to change --- a self-intersecting
design failing loudly instead of simulating silently --- is regression-guarded from the
design file all the way to the thrown, located error. Placement intent in, geometry out,
and a test suite standing over the seam.
